\documentclass[11pt]{article}

\usepackage[margin=1in]{geometry}
\usepackage{amsmath,amssymb,amsfonts,amsthm}
\usepackage{booktabs}
\usepackage{graphicx}
\usepackage{hyperref}
\usepackage{natbib}
\usepackage{enumitem}
\usepackage{xcolor}
\usepackage{multirow}

\newcommand{\Gr}{\mathrm{Gr}}
\newcommand{\R}{\mathbb{R}}
\newcommand{\misscite}[1]{\textcolor{red}{[CITE: #1]}}

\hypersetup{colorlinks=true, linkcolor=blue!60!black, citecolor=blue!60!black, urlcolor=blue!60!black}

\title{When Does Geometry Help Causal Inference?\\
Sheaf Cohomology, Discrete Curvature, and Treatment Heterogeneity\\
in Clinical Epidemiology}

\author{
Elliot Tower
}

\date{July 2026 --- Working Draft}

\begin{document}
\maketitle

\begin{abstract}
Geometric and topological methods are increasingly proposed for causal inference and epidemiological analysis, but it remains unclear when they provide genuine advantages over standard statistical tools.
We address this question through eleven simulation experiments spanning federated data consistency, confound detection, causal graph validation, treatment heterogeneity, and multi-site neuro-epidemiological heterogeneity in multiple sclerosis.

We identify a clear boundary.
On one side, sheaf consistency testing on complete graphs with scalar data reduces algebraically to Cochran's $Q$, and partial correlation achieves perfect confound detection under complete confounding---geometry adds no signal beyond standard tools in these settings.
On the other side, Grassmannian holonomy detects global inconsistency invisible to all pairwise and scalar alternatives (Berry phase $= 1.85$, $p < 0.001$, while maximum pairwise distance remains below the noise threshold), per-edge sheaf $Q$ tests recover heterogeneous DAG structure that a global test misses entirely ($p < 10^{-300}$ vs.\ $p = 0.659$), and $H^1$ effect-modifier classification achieves perfect accuracy with a three-order-of-magnitude separation between transportable and non-transportable effects.

The boundary is structural: geometry helps when data are multivariate (subspace-valued rather than scalar), when consistency constraints span cycles (not just pairs), and when heterogeneity is edge-specific (not global).
We also characterize two failure modes---Ollivier--Ricci curvature fails for edge validation (AUROC $= 0.466$) while Forman--Ricci curvature succeeds (AUROC $= 0.677$), and PCA-based dimensionality reduction destroys treatment effect signal regardless of the upstream CATE estimator---providing practical guidance on method selection.
\end{abstract}

%% ==============================================================
\section{Introduction}
\label{sec:intro}
%% ==============================================================

Clinical epidemiology faces structural problems at several scales.
Federated electronic health record networks aggregate estimates across hospitals that may define phenotypes differently \misscite{federated EHR review}.
Observational biomarker studies confound direct causal effects with associations mediated through shared upstream causes \misscite{confounding in biomarker studies}.
Causal graph learning algorithms produce dense graphs with no reliable post-hoc filter for distinguishing true from spurious edges \misscite{causal discovery review}.
Treatment effect heterogeneity remains difficult to characterize from observational data \misscite{Athey \& Imbens HTE review}.
In multi-site disease studies such as multiple sclerosis, all four problems interact: disease mechanisms vary across patient strata, federated imaging data carry acquisition confounds, and causal DAGs may be heterogeneous across disease phases and treatments \misscite{Lublin 2014 MS phenotype classification}.

Each problem has a geometric aspect.
Federated consistency asks whether local estimates are compatible under known transformations---what sheaf cohomology measures \misscite{Robinson 2014; Curry 2014 sheaves}.
Confound detection asks whether correlations collapse when projected onto the correct conditioning set.
Edge validation asks whether topological features carry information about edge truthfulness.
Treatment subtype recovery asks whether patients cluster in CATE-augmented covariate space.
For multivariate data, the relevant geometry lives on the Grassmannian $\Gr(k,d)$, whose curvature generates Berry phase when local sections are transported around closed loops \misscite{Berry 1984; Simon 1983}.

The question is whether these geometric connections produce practical advantages or merely repackage standard statistics in topological language.
We answer this through eleven experiments organized in two parts.
Part~A (Section~\ref{sec:results-A}) tests geometric methods on standard clinical epidemiology tasks and finds that geometry reduces to existing tools under the assumptions typically satisfied in practice.
Part~B (Section~\ref{sec:results-B}) tests sheaf-cohomological methods on MS-calibrated simulations where those assumptions are violated---multivariate data, cyclic consistency constraints, edge-specific heterogeneity---and finds that geometry detects structure invisible to scalar and pairwise alternatives.
The boundary between the two regimes is the paper's central finding.

%% ==============================================================
\section{Methods}
\label{sec:methods}
%% ==============================================================

\subsection{Sheaf consistency testing}
\label{sec:meth-sheaf}

We model a network of $K$ sites as a graph $G = (V, E)$ where each vertex $v_k$ carries a \emph{stalk}---a local datum that may be scalar (an effect estimate $\hat\beta_k$) or subspace-valued ($V_k \in \Gr(k,d)$).
For each edge $(v_i, v_j)$, a \emph{restriction map} compares the two stalks.
The sheaf Laplacian test statistic for scalar stalks with pairwise-difference restriction maps is
\begin{equation}
    Q = \mathbf{o}^\top \Sigma^{-1} \mathbf{o}, \quad
    \mathbf{o} = D_0 \boldsymbol{\hat\beta}, \quad
    \Sigma = D_0 \operatorname{diag}(\mathrm{se}^2) D_0^\top,
    \label{eq:sheaf-Q}
\end{equation}
where $D_0$ is the signed incidence matrix and $Q \sim \chi^2_{\mathrm{rank}(D_0)}$ under the null of consistent stalks.

For subspace-valued data, consistency is measured through \emph{parallel transport}.
Given sections $V_1, \ldots, V_m \in \Gr(k,d)$ arranged in a cycle, the transport from $V_i$ to $V_{i+1}$ along the geodesic is $T_{i \to i+1} = U V^\top$ where $V_i^\top V_{i+1} = U \Sigma V^\top$ (SVD).
The composed \emph{holonomy} $\Phi = T_{m \to 1} \cdots T_{1 \to 2}$ equals the identity if and only if the sections are globally consistent.
For the linked-column Berry phase construction (two columns of $V_0$ rotating toward shared perpendicular directions with $\pi/2$ phase offset), the holonomy norm is
\begin{equation}\label{eq:berry}
    \|\Phi - I_k\|_F = 2\sqrt{2}\,|\sin(\pi \sin^2 r)|,
\end{equation}
where $r$ is the rotation radius.

\subsection{Per-edge sheaf $Q$ tests}
\label{sec:meth-peredge}

For a causal DAG estimated independently in $S$ strata, we test each edge separately rather than pooling into a global statistic.
For edge $e$ with stratum-specific estimates $\hat\beta_s^{(e)}$ and standard errors $\hat\sigma_s^{(e)}$:
\begin{equation}
    Q^{(e)} = \mathbf{z}^\top \Sigma^{-1} \mathbf{z}, \quad
    z_{ij} = \hat\beta_i^{(e)} - \hat\beta_j^{(e)}, \quad
    \Sigma = D_0 \operatorname{diag}((\hat\sigma_s^{(e)})^2) D_0^\top,
\end{equation}
with $Q^{(e)} \sim \chi^2_{S-1}$ under homogeneity.
Bonferroni correction controls the family-wise error rate.

\subsection{$H^1$ effect-modifier classification}
\label{sec:meth-h1}

For a mechanism--modifier pair, patients are stratified by the modifier into $K$ groups and the causal effect is estimated within each stratum.
The sheaf $Q$ test classifies the pair as \emph{transportable} ($Q$ non-significant, $H^1 \approx 0$) or \emph{stratum-specific} ($Q$ significant, $H^1 \neq 0$).
Sample size per stratum is calibrated to expected interaction strength to avoid over-powered detection of negligible heterogeneity.

\subsection{Confound detection}
\label{sec:meth-confound}

\paragraph{Partial correlation collapse.}
For each biomarker $j$, we compute raw Pearson correlation $r_j = \operatorname{cor}(X_j, Y)$, partial correlation $r_j^{\mathrm{partial}}$ conditioning on all other covariates, and intensive margin correlation $r_j^{\mathrm{int}}$ conditioning on treatment.
Confounded biomarkers' partial correlations should collapse toward zero.

\paragraph{Bracket-norm confound audit.}
For multi-site imaging data, confound leakage is $\Delta = (R^2_{\mathrm{metric}} - R^2_{\mathrm{unique}}) / R^2_{\mathrm{metric}}$, where $R^2_{\mathrm{unique}} = R^2_{\mathrm{both}} - R^2_{\mathrm{acq}}$ is the disease-severity variance unexplained by acquisition parameters.
A metric passes if $\Delta < 0.1$ and the post-correction partial correlation with a biological anchor (sNfL) remains significant.

\subsection{Curvature-based edge validation}
\label{sec:meth-curvature}

We test whether edge-level features discriminate true positive (TP) from false positive (FP) edges in learned causal graphs.
Data are generated from linear and nonlinear structural equation models:
\begin{align}
    \text{Linear:} \quad X_i &= \textstyle\sum_{j \in \mathrm{pa}(i)} W_{ji} X_j + \epsilon_i, \\
    \text{Nonlinear:} \quad X_i &= \textstyle\sum_{j \in \mathrm{pa}(i)} W_{ji} X_j + 0.4 \sin\!\bigl(2 \textstyle\sum_{j} W_{ji} X_j\bigr) + \epsilon_i,
\end{align}
with $W_{ji} \sim \pm\operatorname{Unif}(0.5, 2.0)$ and $\epsilon_i \sim \mathcal{N}(0, 1)$.
Undirected graphs are learned by thresholding the partial correlation matrix.

Six edge features are tested:
Forman--Ricci curvature $\kappa_F(u,v) = 4 - d(u) - d(v)$ \misscite{Forman 2003},
augmented Forman ($\kappa_F + 3|\mathcal{N}(u) \cap \mathcal{N}(v)|$),
Jaccard coefficient,
edge betweenness centrality \misscite{Girvan \& Newman 2002},
partial correlation magnitude $|\hat\rho^{\mathrm{partial}}_{uv}|$,
and average endpoint clustering coefficient.
Ollivier--Ricci curvature (ORC) \misscite{Ollivier 2009; Ni et al.\ 2019} serves as a baseline.

\subsection{Treatment heterogeneity detection}
\label{sec:meth-hte}

We simulate patients with discrete treatment response subtypes and test whether CATE estimation followed by clustering recovers the subtype structure.
Four CATE estimators (KNN T-learner, random forest T-learner, gradient boosting T-learner, random forest S-learner) are crossed with four clustering methods (PCA on CATE-weighted covariates $+$ $K$-means, $K$-means on CATE values, $K$-means on concatenated covariates and CATE, spectral clustering on CATE RBF kernel).
Performance is measured by adjusted Rand index (ARI) against known subtypes.

%% ==============================================================
\section{Results Part A: When Geometry Reduces to Standard Tools}
\label{sec:results-A}
%% ==============================================================

\subsection{Sheaf consistency testing recovers Cochran's $Q$}
\label{sec:res-sheaf}

We simulate $K = 8$ hospital sites with scalar stalks (treatment effect estimates) on a complete graph.
The sheaf test achieves proper calibration (type~I error $= 5.85\%$ at $\alpha = 0.05$, 2{,}000 null replicates) and power increasing monotonically with bias: $38.6\%$ at $\delta = 0.10$, $71.2\%$ at $\delta = 0.15$, $96.8\%$ at $\delta = 0.20$.
The biased site (site~5) is localized to rank~1 among 8 sites in 94\% of simulations.

The sheaf test produces results identical to Cochran's $Q$ at every bias level, structured separation value, and replicate.
On a complete graph with scalar stalks and pairwise-difference restriction maps, the quadratic form in Eq.~\eqref{eq:sheaf-Q} is algebraically equivalent to $Q$.
This equivalence identifies three conditions required for the sheaf framework to add value: heterogeneous stalks, incomplete site graphs, or non-trivial restriction maps.

\subsection{Partial correlation perfectly separates complete confounding}
\label{sec:res-confound}

We simulate $n = 3{,}000$ patients with 10 real and 10 confounded biomarkers (Table~\ref{tab:confound}).
Partial correlation achieves AUROC $= 1.000$ at all sample sizes ($n = 200$ to $5{,}000$) with zero variance.
Confounded biomarkers collapse to partial correlations of $0.001$--$0.033$, while real biomarkers retain $0.61$--$0.76$.

\begin{table}[t]
\centering
\caption{Correlation measures for real vs.\ confounded biomarkers ($n = 5{,}000$).}
\label{tab:confound}
\small
\begin{tabular}{@{}lcc@{}}
\toprule
\textbf{Measure} & \textbf{Real (range)} & \textbf{Confounded (range)} \\
\midrule
Raw Pearson & $0.81$--$0.90$ & $0.51$--$0.71$ \\
Partial correlation & $0.61$--$0.76$ & $0.001$--$0.033$ \\
Intensive margin & $0.80$--$0.89$ & $0.31$--$0.64$ \\
\bottomrule
\end{tabular}
\end{table}

The separation is a property of complete confounding: shared-cause pathways that conditioning fully removes.
Partial confounding (where confounders explain part of the association) and unobserved confounders would reduce the AUROC below $1.0$.

\subsection{ORC fails; Forman--Ricci rescues edge validation}
\label{sec:res-curvature}

Ollivier--Ricci curvature yields AUROC $= 0.466$ for TP vs.\ FP discrimination ($p = 0.020$, below chance).
ORC measures optimal transport between neighborhoods, which depends on global graph topology rather than edge-level statistical evidence.

Forman--Ricci curvature rescues the approach, achieving AUROC $= 0.677$ in the best condition (Table~\ref{tab:curvature}).
The signal is consistent across all six conditions (AUROC $0.634$--$0.677$).
Partial correlation magnitude provides a complementary non-topological signal (AUROC $= 0.657$).

\begin{table}[t]
\centering
\caption{Edge feature AUROC for TP vs.\ FP discrimination (best condition per feature, 60 graphs/condition).}
\label{tab:curvature}
\small
\begin{tabular}{@{}lccc@{}}
\toprule
\textbf{Feature} & \textbf{Best AUROC} & \textbf{Condition} & \textbf{Direction} \\
\midrule
Forman--Ricci & $\mathbf{0.677}$ & linear / loose & TP $>$ FP \\
$|\text{Partial corr}|$ & $0.657$ & nonlinear / loose & TP $>$ FP \\
Betweenness & $0.610$ & linear / loose & TP $>$ FP \\
Avg.\ clustering & $0.568$ & nonlinear / loose & TP $>$ FP \\
Augmented Forman & $0.484$ & --- & wrong direction \\
Jaccard & $0.408$ & --- & wrong direction \\
\midrule
ORC (initial) & $0.466$ & linear & wrong direction \\
\bottomrule
\end{tabular}
\end{table}

The mechanism is degree asymmetry: $\kappa_F = 4 - d(u) - d(v)$ is higher when both endpoints have low degree.
TP edges connect lower-degree nodes on average because real causal relationships produce specific partial correlations, while FP edges arise from indirect paths through higher-degree hub nodes.

\subsection{PCA clustering is the HTE bottleneck}
\label{sec:res-hte}

We simulate $n = 3{,}000$ patients with 3 treatment response subtypes (responders CATE $= +2.0$, anti-responders CATE $= -2.5$, non-responders CATE $= 0$).
The $4 \times 4$ CATE $\times$ clustering grid (Table~\ref{tab:hte}) reveals that clustering method dominates CATE estimator in determining performance.

\begin{table}[t]
\centering
\caption{ARI for HTE subtype recovery ($4 \times 4$ grid, 100 replicates, $n = 3{,}000$).}
\label{tab:hte}
\small
\begin{tabular}{@{}lcccc@{}}
\toprule
 & \textbf{PCA $K$-m.} & \textbf{CATE $K$-m.} & \textbf{CATE+cov} & \textbf{Spectral} \\
\midrule
KNN       & $-0.011$ & $0.179$ & $0.225$ & $0.085$ \\
RF T-learner  & $-0.013$ & $0.238$ & $\mathbf{0.270}$ & $0.097$ \\
GBM T-learner & $-0.016$ & $0.189$ & $0.264$ & $0.055$ \\
RF S-learner  & $-0.014$ & $0.175$ & $0.245$ & $0.066$ \\
\midrule
Oracle        & $0.095$  & $1.000$ & $0.551$ & $1.000$ \\
Naive (raw cov) & \multicolumn{4}{c}{$0.218$} \\
\bottomrule
\end{tabular}
\end{table}

PCA clustering fails universally: all four CATE methods and the oracle (ARI $= 0.095$) produce near-random clusters.
The best estimated combination (RF T-learner + covariate-augmented $K$-means, ARI $= 0.270$) exceeds the naive baseline ($0.218$) by 24\% but falls far short of the oracle ($1.000$).
The bottleneck is CATE estimation noise at $n = 3{,}000$.

\subsection{Part A summary}

These experiments establish a baseline: under the assumptions that typically hold in standard epidemiological settings (scalar data, complete comparisons, complete confounding, moderate sample sizes), geometric methods reduce to or marginally improve upon existing tools.
The sheaf framework recovers Cochran's $Q$.
Partial correlation detection is trivially perfect.
Forman--Ricci curvature provides a genuine topological signal (AUROC $0.677$) but modest in magnitude.
PCA-based dimensionality reduction actively destroys treatment effect signal.

%% ==============================================================
\section{Results Part B: When Geometry Adds Signal}
\label{sec:results-B}
%% ==============================================================

We now test sheaf-cohomological methods on simulations calibrated to multiple sclerosis neurology, where the assumptions violated in Part~A---multivariate data, cyclic constraints, edge-specific heterogeneity---hold by design.

\subsection{Prerequisites}

Cohort contamination rate is $0.9\%$ (below the 2\% threshold), and effect estimates are stable across diagnostic stringency ($p = 0.18$).
IRT-derived theta scores correlate more strongly with sNfL ($r = 0.836$) than raw EDSS ($r = 0.801$), and the latent model reduces scale variability from $0.483$ to $0.082$ (83\% reduction).

\subsection{Grassmannian holonomy detects structure invisible to all alternatives}
\label{sec:res-cocycle}

The cocycle experiment operates on $\Gr(3, 20)$ with $m = 24$ sections forming a closed loop.
We plant a known Berry phase via the linked-column construction at radius $r = 0.5$ and test three conditions simultaneously: (C1) each consecutive pair is close, (C2) the composed transport deviates from identity, and (C3) scalar projections do not detect the structure.

\begin{table}[t]
\centering
\caption{Cocycle obstruction results. The construction achieves all three conditions: pairwise consistent yet globally inconsistent and scalar-blind.}
\label{tab:cocycle}
\small
\begin{tabular}{@{}lcc@{}}
\toprule
\textbf{Metric} & \textbf{Value} & \textbf{Threshold} \\
\midrule
Measured holonomy $\|\Phi - I\|_F$ & $1.851$ & --- \\
Predicted (Eq.~\ref{eq:berry}) & $1.869$ & --- \\
$p$-value vs.\ null & $< 0.001$ & --- \\
\midrule
Max pairwise distance & $0.730$ & $0.767$ (C1 holds) \\
Holonomy norm & $1.851$ & $1.226$ (C2 holds) \\
Scalar projections & non-significant & (C3 holds) \\
\bottomrule
\end{tabular}
\end{table}

The measured holonomy ($1.851$) matches the predicted Berry phase ($1.869$) to within 1\% (Table~\ref{tab:cocycle}).
The construction resolves the C1+C2 paradox: per-step rotation is masked by noise (max pairwise distance $0.730 < 0.767$ threshold), but the composed holonomy accumulates coherently ($1.851 > 1.226$ threshold).
Signal grows as $\sqrt{m}$ while noise accumulates as a random walk, giving SNR $\approx 5$ for $m = 24$.

Three competitor baselines fail to isolate the holonomy.
Maximum pairwise angle cannot distinguish planted from control sections ($0.730$ vs.\ $0.695$, $5\%$ difference).
Random-effects and CKA tests detect subspace spread ($p < 10^{-77}$) but are blind to the cyclic obstruction.
Only composed parallel transport detects the global inconsistency.

\subsection{Per-edge sheaf $Q$ tests recover the two-process MS structure}
\label{sec:res-dag}

A 3-node DAG (inflammation $\leftrightarrow$ degeneration $\to$ disability) is estimated in 8 MS disease strata (Table~\ref{tab:dag}).
The feedback edges vary dramatically across strata: early RRMS shows dominant inflammation$\to$degeneration ($0.404$) with negligible reverse flow ($-0.001$), while PPMS shows the opposite ($-0.008$ and $0.385$).
Disability edges remain stable (variance $\sim\!100\times$ smaller).

\begin{table}[t]
\centering
\caption{Per-edge sheaf $Q$ tests across 8 MS strata. Feedback edges show massive heterogeneity; disability edges are homogeneous.}
\label{tab:dag}
\small
\begin{tabular}{@{}lcccl@{}}
\toprule
\textbf{Edge} & $\boldsymbol{Q}$ & $\boldsymbol{p}$ & \textbf{df} & \textbf{Heterogeneous?} \\
\midrule
infl $\to$ degen & $1806.9$ & $< 10^{-300}$ & 7 & Yes \\
degen $\to$ infl & $1549.6$ & $< 10^{-300}$ & 7 & Yes \\
infl $\to$ disab & $7.05$ & $0.423$ & 7 & No \\
degen $\to$ disab & $9.98$ & $0.189$ & 7 & No \\
\bottomrule
\end{tabular}
\end{table}

Per-edge testing is essential: the global Frobenius test (pooling all edges) produces $p = 0.659$ because near-constant disability edges dilute the signal from heterogeneous feedback edges.
Per-edge tests with Bonferroni correction resolve this, recovering exactly the predicted two-process disease structure.

Treatment effects align with known mechanisms: BTK inhibition suppresses infl$\to$degen from $0.404$ to $0.002$; siponimod preserves degen$\to$disab ($0.386$), consistent with relapse-independent progression \misscite{Kappos 2018 siponimod}.

\subsection{$H^1$ classification achieves perfect accuracy with clean separation}
\label{sec:res-h1}

Seven mechanism--modifier pairs spanning three edge types are classified as transportable or stratum-specific (Table~\ref{tab:h1}).

\begin{table}[t]
\centering
\caption{$H^1$ effect-modifier classification. All 7 pairs correctly classified with a three-order-of-magnitude gap.}
\label{tab:h1}
\small
\begin{tabular}{@{}llccrl@{}}
\toprule
\textbf{Pair} & \textbf{Expected} & $\boldsymbol{\gamma}$ & $\boldsymbol{Q}$ & $\boldsymbol{p}$ & \textbf{Correct} \\
\midrule
HLA $\times$ EBV & transport & $0.10$ & $3.30$ & $0.348$ & Yes \\
EBV necessity & transport & $0.05$ & $6.47$ & $0.091$ & Yes \\
Vitamin D & transport & $0.04$ & $6.97$ & $0.073$ & Yes \\
\midrule
Sex $\times$ course & non-transp & $0.50$ & $2434$ & $< 10^{-300}$ & Yes \\
Genetics $\times$ OCB & non-transp & $0.60$ & $3588$ & $< 10^{-300}$ & Yes \\
Age $\times$ anti-CD20 & non-transp & $0.40$ & $1705$ & $< 10^{-300}$ & Yes \\
Phenotype $\times$ GM & non-transp & $0.45$ & $1828$ & $< 10^{-300}$ & Yes \\
\bottomrule
\end{tabular}
\end{table}

The separation is bimodal: transport pairs have $Q < 7$ ($p > 0.07$) and non-transport pairs have $Q > 1700$ ($p < 10^{-300}$), spanning three orders of magnitude.
All transport pairs are exposure$\to$mechanism edges with weak interaction ($\gamma \leq 0.10$), matching the domain expectation that causal risk factors operate through mechanisms consistent across patient strata.

\subsection{Bracket-norm confound audit confirms all four MS imaging metrics}
\label{sec:res-bracket}

Four MS imaging biomarkers (iron rim QSM, deep GM atrophy, cortical lesion count, cervical cord CSA) are tested across $n = 1{,}500$ patients at 8 sites.
All metrics pass: confound leakage $\Delta < 0.05$ for three metrics and $\Delta = -0.19$ (suppressor effect) for iron rim QSM.
Post-correction partial correlations with sNfL remain significant ($p < 10^{-55}$ for all).

%% ==============================================================
\section{Discussion}
\label{sec:discussion}
%% ==============================================================

\subsection{The boundary}

Parts A and B together identify a structural boundary for when geometric methods add signal.

\paragraph{Geometry helps when data are subspace-valued.}
Scalar stalks on a complete graph reduce the sheaf test to Cochran's $Q$ (Section~\ref{sec:res-sheaf}).
Subspace-valued stalks on the Grassmannian produce Berry phase---holonomy that accumulates coherently around cycles while noise accumulates as a random walk (Section~\ref{sec:res-cocycle}).
The transition from scalar to multivariate data is the primary determinant of whether geometric structure carries information.

\paragraph{Geometry helps when constraints span cycles.}
Pairwise comparisons cannot detect the cocycle obstruction: maximum pairwise distance differs by only 5\% between planted and control conditions.
The holonomy, by contrast, separates them completely ($1.851$ vs.\ $0.169$).
This gap arises from the $\sqrt{m}$ signal accumulation of coherent transport around an $m$-step cycle, which has no pairwise analogue.

\paragraph{Geometry helps when heterogeneity is edge-specific.}
Global heterogeneity tests average high-variance edges with near-constant ones.
In the MS DAG simulation, this averaging produces a false negative ($p = 0.659$) for the global test, while per-edge tests correctly identify the heterogeneous feedback edges ($p < 10^{-300}$).

\subsection{Failure modes}

Two negative results provide practical method-selection guidance.

\paragraph{Curvature type matters.}
ORC measures optimal transport between neighborhoods (sensitive to global topology), while Forman--Ricci measures degree deficit (sensitive to local connectivity).
For edge validation in learned causal graphs, the relevant structural property is degree asymmetry between TP and FP edges, which Forman--Ricci captures directly and ORC does not.
The lesson generalizes: the ``right'' curvature depends on which structural property discriminates the classes of interest.

\paragraph{PCA destroys treatment effect signal.}
PCA on CATE-weighted covariates produces near-random clusters regardless of CATE estimator quality---even the oracle achieves only ARI $= 0.095$.
The multiplicative interaction between noisy CATE estimates and covariate values introduces variance that PCA's linear projection amplifies rather than filters.
Direct clustering on CATE estimates or CATE-augmented covariates consistently outperforms the PCA pathway.

\subsection{Limitations}

All experiments use simulated data designed to test specific methodological claims.
The confound audit uses complete confounding; the SEM simulations use random DAGs without domain structure; the HTE simulation uses well-separated subtypes; the MS simulations plant known heterogeneity patterns.
Real clinical data present partial confounding, domain-structured graphs, continuous treatment effect variation, measurement error, and missing data.

The cocycle dose-response (ARM~4) shows substantial point-to-point noise at small radii, suggesting that larger $m$ or repeated trials would be needed for smooth dose-response in practice.
The $H^1$ classifier's scale-invariance stress test did not fully pass: quantile normalization reduced scale variability from $0.517$ to $0.346$ but did not reach the $0.30$ target.

\subsection{Neuro-epidemiological perspective}

The MS results illustrate a broader principle for neuro-epidemiological study design.
The two-process model (inflammation$\leftrightarrow$degeneration$\to$disability) generates heterogeneity structure that is invisible to global tests but recoverable by per-edge sheaf analysis.
This suggests that multi-site MS cohort studies should test DAG edges individually rather than pooling across the entire causal structure, and that transportability assessment (which effects generalize across patient strata) should be a standard step before pooling multi-site estimates.

The cocycle obstruction method extends naturally to multi-site imaging studies where each site's biomarker covariance matrix defines a Grassmannian section.
Sites with consistent measurement produce zero holonomy; sites with acquisition-dependent distortions produce non-zero holonomy that localizes the inconsistency.
This is precisely the setting where Cochran's $Q$ cannot be applied (the data are multivariate, the site graph may be incomplete, and the restriction maps are non-trivial).

%% ==============================================================
\section{Next Experiments}
\label{sec:next}
%% ==============================================================

\begin{enumerate}[nosep]
    \item \textbf{Multi-phenotype sheaf}: Sites measure different phenotype features (labs, imaging, codes) with overlapping coverage. Test whether sheaf cohomology detects inconsistencies that Cochran's $Q$ cannot express.

    \item \textbf{Partial confounding spectrum}: Grade confound strength from $0\%$ to $100\%$ and measure detection AUROC at each level. Include unobserved confounders.

    \item \textbf{Combined topological-statistical edge classifier}: Train on the Forman--Ricci $+$ partial correlation magnitude feature set. Test on KEGG/STRING/BioGRID graph structures.

    \item \textbf{Realistic HTE}: Smaller CATE differences ($0.3$--$0.5$ SD), more subtypes, continuous boundaries, time-varying effects, larger cohorts.

    \item \textbf{Real-data imaging cocycle}: Apply Grassmannian holonomy to multi-site MRI covariance matrices from a real MS cohort.

    \item \textbf{Scale-free $H^1$ classifier}: Develop alternatives to OLS-based sheaf $Q$ that are invariant to outcome scale choice.
\end{enumerate}

\paragraph{Code and data availability.}
All simulation code and results are available at \textcolor{red}{[REPO URL]}.

\bibliographystyle{unsrtnat}
% \bibliography{references}

\end{document}
